% Generated draft: requires mathematical and source review.
\documentclass[11pt]{article}
\usepackage[margin=25mm]{geometry}
\usepackage{fontspec}
\setmainfont{Latin Modern Roman}
\usepackage{amsmath,amssymb,mathtools,unicode-math}
\setmathfont{Latin Modern Math}
\usepackage{enumitem,needspace,xurl,hyperref}
\setlength{\parindent}{0pt}
\setlength{\parskip}{5pt}
\setlength{\emergencystretch}{3em}
\newcommand{\N}{\mathbb{N}}
\newcommand{\Z}{\mathbb{Z}}
\newcommand{\Q}{\mathbb{Q}}
\newcommand{\R}{\mathbb{R}}
\newcommand{\C}{\mathbb{C}}
\newcommand{\F}{\mathbb{F}}
\newcommand{\A}{\mathbb{A}}
\newcommand{\PP}{\mathbb{P}}
\newcommand{\E}{\mathbb{E}}
\newcommand{\eps}{\varepsilon}
\newcommand{\dd}{\,\mathrm{d}}
\begin{document}
\section*{UnsolvedMath ID: 20002884}
\textbf{Source problem code:} catalogue entry
\subsection{Definitions and mathematical statement}
John D'Angelo and Dror Varolin ask for the power radical of polynomial squared norms. Let $\mathcal P$ be Hermitian polynomial kernels $R(z,\bar w)=\overline{R(w,\bar z)}$, and let $\mathcal S$ consist of kernels $\sum_jp_j(z)\overline{p_j(w)}$ with holomorphic polynomials $p_j$. The target is to characterize
\[
 \{R\in\mathcal P:\exists N\ge1,\ R^N\in\mathcal S\}.
\]
Equality on the diagonal polarizes to kernel equality. Positive exponents are intended; the zero kernel is a trivial case. The original OCR loses several conjugation bars and superscripts.

The source records two exact criteria only for $N=1$. First, $R\in\mathcal S$ precisely when every finite evaluation matrix $(R(z_i,\bar z_j))_{i,j=1}^k$ is positive semidefinite, for all $k$ and all point choices. Second, writing $R=\sum_{\alpha,\beta}c_{\alpha\beta}z^\alpha\bar w^\beta$, it gives positivity for every finite-size commuting tuple $A=(A_1,\ldots,A_n)$:
\[
 R[A,A^*]=\sum_{\alpha,\beta}c_{\alpha\beta}(A^*)^\beta A^\alpha\ge0.
\]
Adjoints are placed first; the $A_i$ commute with one another but need not be normal or commute with adjoints. Neither $N=1$ test by itself classifies roots of squared norms.

The aid reports that a nonzero $R$ admitting a power has constant weak diagonal sign off its complex-algebraic zero set. Choose $Q=R$ or $-R$ with nonnegative diagonal. It satisfies the polarized two-point inequality
\[
 |Q(z,\bar w)|^2\le Q(z,\bar z)Q(w,\bar w).
\]
The source already notes that this condition can hold without being a squared norm. The aid adds equality rigidity: if $Q(a,\bar a),Q(b,\bar b)>0$ and equality holds for $(a,b)$, then
\[
 Q(u,\bar b)=\frac{Q(a,\bar b)}{Q(a,\bar a)}Q(u,\bar a)
 \quad\text{as a polynomial identity in }u.
\]
Thus a positive-power representation forces proportional full polarized rows at every such equality pair, not merely the two-point numerical equality.

Combining this with the cited low-rank theory, the aid reports a complete special case: if the Hermitian coefficient matrix has rank at most two, some positive power belongs to $\mathcal S$ exactly when $R\in\mathcal S$ or $-R\in\mathcal S$. The negative case is permitted by even powers.

It also treats, in $\C^2$, the family
\[
 r_{k,t}(z,\bar w)=(z_1\bar w_1+z_2\bar w_2)^{2k}
 -t(z_1z_2\bar w_1\bar w_2)^k.
\]
For every integer $k\ge2$, the aid states that $t=4^k/2$ is the sharp Cauchy--Schwarz boundary member, satisfies that inequality, and has no squared-norm power. The equality pair $a=(1,1)$, $b=(1,-1)$ has positive diagonals and equal absolute cross value, but its two polarized rows are not proportional, supplying the stated obstruction. The restriction $k\ge2$ is essential. These results do not give the general necessary-and-sufficient characterization, which the aid describes as open in its checked literature.
\subsection{Short English statement}
The source's first-power tests are kept distinct from its power-radical question. The aid supplies an additional equality-row obstruction, a rank-at-most-two classification, and Cauchy–Schwarz boundary examples with no squared-norm power. It leaves the general characterization unresolved.
\subsection{Sources}
\begin{itemize}
\item Local UnsolvedMath registry, permanent ID 20002884, source code catalogue entry; created and updated 14 August 2026. Canonical \texttt{aim-several-complex-variables-notes.json}, zero-based \texttt{notes[41]}, tag \texttt{problem}, item 2 proposed by John D'Angelo and Dror Varolin.
\item AIM ``Complexity of mappings in CR geometry'': \url{https://aimath.org/WWN/complexitycrmap/complexitycrmap.pdf}. The record is marked exact although the supplied verification reports lost bars/superscripts and a page footer; it says the PDF and references were inspected. No inspection is claimed here. Original remarks/literature fields are empty.
\item The source cites [CD] and [DV] without their expanded bibliographic entries. Research attempt 1: \texttt{partial\_result}, problem status at run open, rigor label \texttt{proved}, novelty confidence low. No complete authoritative references or theorem hypotheses are supplied beyond those labels and the D'Angelo--Varolin attribution.
\item Ulam AI / UnsolvedMath Contributors, CC BY 4.0. Machine-generated research aid, not peer reviewed; independently verify before citation.
\end{itemize}
\end{document}
